% TOP_PROBLEM: {"id": 30007373, "title": "Spherical perceptron pressure with orthogonally invariant disorder", "category_id": 19, "category": {"id": 19, "name": "probability", "display_name": "Probability", "description": "Problems involving probability theory, stochastic processes, and random structures.", "slug": "probability", "order_index": 19, "created_at": "2026-07-31T15:26:25.671Z"}, "status": "open"}
% FIRST_POSTED: 2026-10-06
% ATTEMPT_STATUS: solved
% OPENAI_CLAIM: {"schema_version":1,"source_repo":"https://github.com/openai/math","source_commit":"adc7f1241b42e322a6451854ab7e4b4c146bf78a","release_date":"2026-10-06","match":"full","resolution":"proved","independently_reviewed":false,"second_pass":"agreed","families":[{"family":"222","title":"Perceptron free energies and microscopic jamming exponents","subject":"Probability and statistical mechanics","manuscripts":[{"title":"The spherical perceptron with bi-orthogonally invariant disorder","pdf_path":"preprints/The-spherical-perceptron-with-bi-orthogonally-invariant-disorder-September-24-2026/The-spherical-perceptron-with-bi-orthogonally-invariant-disorder-September-24-2026.pdf","readme_path":"preprints/The-spherical-perceptron-with-bi-orthogonally-invariant-disorder-September-24-2026/README.md","date":"2026-09-24","theorem_ref":"thm:main (paper.tex:92); prop:aspect (paper.tex:1102)","statement_sources":[{"path":"preprints/The-spherical-perceptron-with-bi-orthogonally-invariant-disorder-September-24-2026/build/paper.tex","line":92,"label":"thm:main"},{"path":"preprints/The-spherical-perceptron-with-bi-orthogonally-invariant-disorder-September-24-2026/build/paper.tex","line":1102,"label":"prop:aspect"}]}],"lean":{"doc_path":"lean/docs/222.md","comparators":[],"declarations":[],"covers_main_theorem":false},"reasoning_trace":null}]}
\documentclass[11pt]{article}
\usepackage[margin=25mm]{geometry}
\usepackage{fontspec}
\setmainfont{Latin Modern Roman}
\usepackage{amsmath,amssymb,mathtools}
\usepackage{unicode-math}
\setmathfont{Latin Modern Math}
\usepackage{enumitem}
\usepackage{xurl,hyperref,bookmark}
\hypersetup{colorlinks=true,urlcolor=blue}
\setcounter{secnumdepth}{0}
\setlength{\parindent}{0pt}
\setlength{\parskip}{5pt}
\setlength{\emergencystretch}{3em}
\begin{document}
\section*{Spherical perceptron pressure with orthogonally invariant disorder}

\subsection{OpenAI's claim}
OpenAI presents a mathematical result relevant to this question in the selected manuscript. A spherical perceptron with bi-orthogonally invariant disorder is claimed to have an explicit pressure formula for bounded continuous activation and compact singular-value limits without outliers.

\subsection{Mathematical statement of the claimed result}
\paragraph{Result 1.} Let $M/N\to\alpha>0$, $r_N=\min(M,N)$, and let the law of a real $M\times N$ matrix $A_N$ be invariant under left and right orthogonal multiplication. Suppose its singular values $s_{N,r}$ satisfy, in probability,
\[
 r_N^{-1}\sum_{r=1}^{r_N}\delta_{s_{N,r}}\Longrightarrow\nu,\qquad
 \max_{r\le r_N}\operatorname{dist}(s_{N,r},\operatorname{supp}\nu)\to0,
\]
for a deterministic compactly supported law $\nu$ on $[0,\infty)$. For bounded continuous $f$, with $\sigma_N$ the uniform probability on $\|x\|^2=N$, put
\[
 F_N=N^{-1}\log\int e^{\sum_{r=1}^Mf((A_Nx)_r)}\,d\sigma_N(x).
\]
Then $F_N$ converges in probability and mean to a deterministic $F_{\alpha,\nu}(f)$. Its finite spectral-block formula is
\[
 \sup_{\rho_i>0,\ \sum_i\rho_i=1}
 \left\{\frac12\sum_iv_i\log(\rho_i/v_i)+
 \alpha P_p^f\bigl((s_j^2\rho_j/\alpha)_{j=1}^h\bigr)\right\}.
\]
Here positive singular-value blocks have $n_j/M\to p_j>0$, $v_j=\alpha p_j$, and $\sum_jp_j=\min(1,1/\alpha)$; when $\alpha<1$ an input null block has $v_0=1-\alpha$. The angular functional $P_p^f(a)$ is the deterministic limit, as $M\to\infty$, of
\[
 M^{-1}\log\int\exp\left(\sum_{r=1}^Mf\left(\sum_jy_{j,r}\right)\right)
 \prod_jd\sigma_{E_j,Ma_j}(y_j),
\]
where mutually orthogonal Haar subspaces $E_j\subset\mathbb R^M$ have dimension ratios $p_j$, and $\sigma_{E_j,Ma_j}$ is uniform probability on their spheres of squared radius $Ma_j$ (a point mass for $a_j=0$). Taking increasingly fine operator-norm spectral approximations gives $F_{\alpha,\nu}(f)$. Blocks of vanishing proportion are omitted.

\subsection{Outline of the argument}
Marked probability cascades encode the overlap hierarchy while orthogonal invariance controls coordinate cavity increments. Conditional change-of-measure calculations and spectral data determine the disorder contribution; smooth compactly supported activations are handled first and then approximation removes that regularity.

\subsection{Scope relative to this problem}
The main theorem gives the requested deterministic pressure limit and spectral-block variational formula for bounded continuous activation and compact singular spectra without outliers. The aspect-ratio extension permits arbitrary deterministic $M/N\to\alpha$, removing the main statement's floor convention.

\subsection{Formal verification}
For family 222, the spherical Gaussian perceptron pressure formula is included. The Ising companion entry only proves finite variational value for bounded continuous log-potentials; it neither gives its pressure convergence nor the bounded-Borel extension. A spherical-field supplement supplies a dual approximation formula. No Comparator is assigned to the selected manuscript here; the family page selects other papers or supporting statements. This site has not executed the Lean code. The scope descriptions do not certify the remaining manuscript arguments.

\subsection{Status caveat}
OpenAI attributes these manuscripts to work by an internal model. Its release cautions that arguments without formalization can contain mistakes. This catalogue has not independently checked the proofs; the labels report the mathematical claims attributed to OpenAI.

\subsection{Sources}
\begin{itemize}
\item OpenAI, \emph{The spherical perceptron with bi-orthogonally invariant disorder}, 2026-09-24; citation key \texttt{OAI:The-spherical-perceptron-with-bi-orthogonally-invariant-disorder-September-24-2026}. \href{https://github.com/openai/math/blob/adc7f1241b42e322a6451854ab7e4b4c146bf78a/preprints/The-spherical-perceptron-with-bi-orthogonally-invariant-disorder-September-24-2026/The-spherical-perceptron-with-bi-orthogonally-invariant-disorder-September-24-2026.pdf}{PDF (GitHub)}; \href{https://github.com/openai/math/raw/adc7f1241b42e322a6451854ab7e4b4c146bf78a/preprints/The-spherical-perceptron-with-bi-orthogonally-invariant-disorder-September-24-2026/The-spherical-perceptron-with-bi-orthogonally-invariant-disorder-September-24-2026.pdf}{PDF (direct)}; \href{https://github.com/openai/math/blob/adc7f1241b42e322a6451854ab7e4b4c146bf78a/preprints/The-spherical-perceptron-with-bi-orthogonally-invariant-disorder-September-24-2026/README.md}{Manuscript page and citation}; \href{https://github.com/openai/math/blob/main/preprints/The-spherical-perceptron-with-bi-orthogonally-invariant-disorder-September-24-2026/The-spherical-perceptron-with-bi-orthogonally-invariant-disorder-September-24-2026.pdf}{Latest version}; Statement source: \href{https://github.com/openai/math/blob/adc7f1241b42e322a6451854ab7e4b4c146bf78a/preprints/The-spherical-perceptron-with-bi-orthogonally-invariant-disorder-September-24-2026/build/paper.tex\#L92}{thm:main}; Statement source: \href{https://github.com/openai/math/blob/adc7f1241b42e322a6451854ab7e4b4c146bf78a/preprints/The-spherical-perceptron-with-bi-orthogonally-invariant-disorder-September-24-2026/build/paper.tex\#L1102}{prop:aspect}.
\item Family 222: \href{https://github.com/openai/math/blob/adc7f1241b42e322a6451854ab7e4b4c146bf78a/lean/docs/222.md}{Lean scope}.
\item \href{https://github.com/openai/math/blob/adc7f1241b42e322a6451854ab7e4b4c146bf78a/CONTENTS.md}{OpenAI catalogue entry}.
\end{itemize}
\end{document}
